% =============================================================================
% The LQ Digital Mode Family — Reference Implementation (lq_lib)
%
% Author:  Luis Quesada (HB9IPH)
% Web:     https://luisquesada.com
% Portal:  https://lquesada.github.io/lq_lib/
% GitHub:  https://github.com/lquesada/lq_lib
% App:     qFT8 — Portable Amateur Radio for Android (https://qft8.com)
%
% License: MIT License (https://github.com/lquesada/lq_lib/blob/main/LICENSE)
%
% Copyright (c) 2026 Luis Quesada (HB9IPH)
%
% Permission is hereby granted, free of charge, to any person obtaining a copy
% of this software and associated documentation files (the "Software"), to deal
% in the Software without restriction, including without limitation the rights
% to use, copy, modify, merge, publish, distribute, sublicense, and/or sell
% copies of the Software, and to permit persons to whom the Software is
% furnished to do so, subject to the following conditions:
%
% The above copyright notice and this permission notice shall be included in
% all copies or substantial portions of the Software.
%
% THE SOFTWARE IS PROVIDED "AS IS", WITHOUT WARRANTY OF ANY KIND, EXPRESS OR
% IMPLIED, INCLUDING BUT NOT LIMITED TO THE WARRANTIES OF MERCHANTABILITY,
% FITNESS FOR A PARTICULAR PURPOSE AND NONINFRINGEMENT. IN NO EVENT SHALL THE
% AUTHORS OR COPYRIGHT HOLDERS BE LIABLE FOR ANY CLAIM, DAMAGES OR OTHER
% LIABILITY, WHETHER IN AN ACTION OF CONTRACT, TORT OR OTHERWISE, ARISING FROM,
% OUT OF OR IN CONNECTION WITH THE SOFTWARE OR THE USE OR OTHER DEALINGS IN THE
% SOFTWARE.
% =============================================================================

% ===========================================================================
%  Section III: Message Architecture (ALLOCATION: PAGE 3 & PAGE 4)
%  Section text MUST fit entirely on Page 3.
%  Dedicated full-page Tables 2, 3, 4 (table*) MUST occupy Page 4.
%  Terminates with \newpage.
% ===========================================================================

\section{Message Architecture}
\label{sec:msg_arch}
\label{sec:typedetails}
\label{sec:huffman}

To maximize channel efficiency within the 77-bit physical frame budget, \LQ{} replaces fixed-width bit fields with an entropy-matched variable-length message architecture.

\subsection{Design Motivation \& Prefix-Free Framing}
\label{sec:huffman:motivation}

Under a strict 77-bit physical payload budget, conventional fixed-width message headers impose an unnecessary penalty on high-entropy transmissions. In a standard \msgtype{CALL} reply, encoding two 28-bit standard callsigns, a 15-bit Maidenhead grid locator, and a 5-bit measured SNR report requires 76 information bits:
\begin{equation}
  \underbrace{1\text{b}}_{\text{Prefix}} + \underbrace{28\text{b}}_{\text{Target}} + \underbrace{28\text{b}}_{\text{Caller}} + \underbrace{15\text{b}}_{\text{Locator}} + \underbrace{5\text{b}}_{\text{SNR}} = \mathbf{77\text{ bits}}.
\label{eq:call_budget}
\end{equation}
Uniform 3-bit or 4-bit type headers exceed the remaining 1-bit margin, preventing simultaneous transmission of locator and report in conventional protocols.

\LQ{} resolves this constraint by synthesizing variable-length prefix coding~\cite{huffman,shannon} with physical-layer bit budgets. Assigning a \textbf{1-bit} codeword (\texttt{1}) to the dominant \msgtype{CALL std (nosuf)} format preserves 76 bits for payload data, combining station identification, locator delivery, and signal report measurement into a single frame. Simpler messages (such as standard CQ, requiring 43 payload bits) readily accommodate longer prefix codes.

Formally, the prefix code $\mathcal{C} = \{c_1, c_2, \ldots, c_k\}$ is \emph{prefix-free}:
\begin{equation}
  \forall\, i \neq j: \quad c_i \text{ is not a prefix of } c_j.
\end{equation}
The receiver traverses a binary prefix tree bit-by-bit from the most-significant bit until reaching a terminal leaf node, identifying the message type without length fields or delimiters. Although the mathematical binary prefix tree across all 16 slots satisfies $\sum_{i=1}^{16} 2^{-\text{len}_i} = 1.0$, the active operational codebook is intentionally prefix-free without being Kraft-complete because Types~14--16 are explicitly reserved for future backward-compatible protocol expansions rather than active operational message formats.

\subsection{Operational Categories \& Message Catalogue}
\label{sec:msg:catalogue}

Stations select transmission formats according to the operational categories defined in Table~\ref{tab:categories_of_operation}. Table~\ref{tab:huffman} catalogues all 16 message slots with their prefix codewords, payload layouts, and bit budgets. Frame payloads $\le 77$ bits are zero-padded to the 77-bit link-layer boundary. Payloads are serialized in big-endian bit order (MSB first) into a 10-byte buffer: bit 0 corresponds to \texttt{0x80} of Byte 0, bit 7 to \texttt{0x01} of Byte 0, and bits 77--79 (the three least significant bits of Byte 9) are set to zero.

\textbf{Field Designations in Table~\ref{tab:huffman}:}
\begin{itemize}[leftmargin=1.4em, itemsep=1.5pt plus 1.0pt minus 0.5pt, topsep=2pt plus 1.0pt minus 0.5pt, parsep=0pt]
  \item \textbf{Canonical Short Identifiers:} Standardized mnemonics: \textbf{CQ}, \textbf{CALL}, \textbf{RPT73} (for \msgtype{REPORT+73}), \textbf{73}, \textbf{RPT73M} (for \msgtype{MULTI-REPORT+73}), and \textbf{73M} (for \msgtype{MULTI-73}).
  \item \textbf{Prefix \& Standard Callsign:} ``Code'' denotes prefix length; ``call:std'' is a 28-bit standard callsign (Section~\ref{sec:std_call}).
  \item \textbf{Non-Standard Callsigns \& Suffixes:} ``call:$n$'' is an $n$-character non-standard callsign packed via Base-38 (Section~\ref{sec:nonstd_call}); ``Suf'' encodes 1-bit portable modifiers (\texttt{0}=\emph{none}, \texttt{1}=\texttt{/P}). Type~5 omits suffixes to pack two standard callsigns, locator, and SNR into 77 bits ($1+28+28+15+5=77\text{b}$).
  \item \textbf{Callsign Hashes:} ``Hash'' is a 24-bit CRC-24/Q hash ($2^{24} \approx 16.78\times 10^6$ bins); ``hash$^\dagger$'' is the 16-bit DX hash (Section~\ref{sec:hash}).
  \item \textbf{Multi-Station Fields:} Type~11 Fox pileup sequences $[\text{Prefix}_3 \mid \text{DX}_{16} \mid \text{Tgt1}_{24} \mid \text{SNR1}_5 \mid \text{Tgt2}_{24} \mid \text{SNR2}_5]$ (77 bits); Type~12 formats $[\text{Prefix}_8 \mid \text{DX}_{16} \mid \text{Tgt1}_{24} \mid \text{Tgt2}_{24} \mid 00000_2]$ (77 bits, 5 zero-padded bits).
  \item \textbf{Modifiers, Locators \& SNR:} ``Mod'' is a 20-bit modifier (Section~\ref{sec:modifier}); ``loc'' is a 15-bit locator (Section~\ref{sec:locator}); ``SNR'' is a 5-bit report (Section~\ref{sec:rst}).
\end{itemize}

\begin{table*}[p]
\centering

% --- Table 2: Operational Categories ---
\small
\setlength{\tabcolsep}{4pt}
\renewcommand{\arraystretch}{0.95}
\caption{Operational categories and corresponding 4-step message sequences.}
\label{tab:categories_of_operation}
\begin{tabular*}{\textwidth}{@{\extracolsep{\fill}}clcccc@{}}
\toprule
\textbf{Cat.} & \textbf{Callsign / Mode Configuration} & \textbf{Step 1: CQ} & \textbf{Step 2: CALL} & \textbf{Step 3: REPORT+73} & \textbf{Step 4: 73} \\
\midrule
\textbf{A} & Standard (no suffixes) & Type 1 & Type 5 & Type 8 & Type 9 \\
\textbf{B} & Standard + std suffix (\texttt{/P}) & Type 1 & Type 6 & Type 8 & Type 9 \\
\textbf{C} & Non-standard ($\le 9$ chars) + optional suffix & Type 2 / 3 & Type 7$^*$ / Type 6 & Type 11 & Type 10 / 12 \\
\textbf{D} & Long compound ($10$--$13$ chars) + optional suffix & Type 4 & ---$^\ddagger$ & Type 11 & Type 12 \\
\textbf{E} & Multi-station reply & --- & --- & Type 11 (Fox) & \makecell{Type 12 (Fox) /\\ Type 9, 10 (Hounds)} \\
\bottomrule
\multicolumn{6}{@{}p{\textwidth}@{}}{\vspace{2pt}\footnotesize $^*$In Step~2, Category~C callers transmit Type~7 (clear-text callsign, suffix, SNR to target 20-bit hash); standard callers answering a Category~C CQ transmit Type~6 (locator, SNR to target 24-bit hash). \quad $^\ddagger$Category~D compound stations ($10$--$13$ chars) initiate contact via CQ (Type~4) or free text (Type~13).}
\end{tabular*}

\vspace{6pt plus 1fill}

% --- Table 3: Complete Message Catalogue ---
\caption{Complete \LQ{} message-type catalogue.}
\label{tab:huffman}
\small
\setlength{\tabcolsep}{2.5pt}
\renewcommand{\arraystretch}{0.94}
\begin{tabular*}{\textwidth}{@{\extracolsep{\fill}}lclccccccccr@{}}
\toprule
\textbf{Huffman Code} & \textbf{ID} & \textbf{Type} & \multicolumn{8}{c}{\textbf{Bits (Payload)}} & \textbf{Total} \\
\cmidrule(lr){4-11}
& & & \textbf{Code} & \textbf{Target} & \textbf{Suf} & \textbf{Caller} & \textbf{Suf} & \textbf{Modifier} & \textbf{Locator} & \textbf{SNR} & \\
\midrule
\bitfield{\szero\szero\szero\szero\szero\szero\szero\szero\szero\szero}         & 1  & CQ std+suf                      & 10 & ---           & ---   & 28 call:std   & 1 suf & 20 mod & 15 loc & ---   & 74 \\
\bitfield{\szero\szero\szero\szero\szero\szero\szero\szero\szero1}              & 2  & CQ non-std~1                    & 10 & ---           & ---   & 48 call:9     & 1 suf & ---    & 15 loc & ---   & 74 \\
\bitfield{\szero\szero\szero\szero\szero\szero\szero1}                          & 3  & CQ non-std~2                    & 8  & ---           & ---   & 48 call:9     & 1 suf & 20 mod & ---    & ---   & 77 \\
\bitfield{\szero\szero\szero\szero\szero\szero1}                                  & 4  & CQ non-std~3                    & 7  & ---           & ---   & 69 call:13    & 1 suf & ---    & ---    & ---   & 77 \\
\bitfield{1}                                                                  & 5  & CALL std (nosuf)                & 1  & 28 call:std   & ---$^*$ & 28 call:std & ---$^*$ & ---  & 15 loc & 5 SNR & 77 \\
\bitfield{\szero1\szero\szero}                                                & 6  & CALL std+suf                    & 4  & 24 hash       & ---   & 28 call:std   & 1 suf & ---    & 15 loc & 5 SNR & 77 \\
\bitfield{\szero\szero1}                                                       & 7  & CALL non-std                    & 3  & 20 hash       & ---   & 48 call:9     & 1 suf & ---    & ---    & 5 SNR & 77 \\
\bitfield{\szero\szero\szero\szero\szero\szero\szero\szero1\szero}              & 8  & REPORT+73 std+suf (RPT73)       & 10 & 28 call:std   & 1 suf & 28 call:std   & 1 suf & ---    & ---    & 5 SNR & 73 \\
\bitfield{\szero\szero\szero\szero\szero\szero\szero\szero11}                  & 9  & 73 std+suf                      & 10 & 28 call:std   & 1 suf & 28 call:std   & 1 suf & ---    & ---    & ---   & 68 \\
\bitfield{\szero\szero\szero1}                                                & 10 & 73 non-std                      & 4  & 24 hash       & ---   & 48 call:9     & 1 suf & ---    & ---    & ---   & 77 \\
\bitfield{\szero11}                                                           & 11 & MULTI-REPORT+73 (RPT73M)        & 3  & \multicolumn{7}{c}{16 hash$^\dagger$ \quad+\quad 2 $\times$ (24 hash + 5 SNR)} & 77 \\
\bitfield{\szero\szero\szero\szero\szero111}                                  & 12 & MULTI-73 (73M)                  & 8  & \multicolumn{7}{c}{16 hash$^\dagger$ \quad+\quad 2 $\times$ 24 hash}            & 72 \\
\bitfield{\szero1\szero1}                                                     & 13 & FREE TEXT                       & 4  & \multicolumn{7}{c}{73 text (Varicode)} & 77 \\
\bitfield{\szero\szero\szero\szero1}                                           & 14 & \emph{(reserved A)}             & 5  & \multicolumn{7}{c}{72 payload} & 77 \\
\bitfield{\szero\szero\szero\szero\szero1\szero}                              & 15 & \emph{(reserved B)}             & 7  & \multicolumn{7}{c}{70 payload} & 77 \\
\bitfield{\szero\szero\szero\szero\szero11\szero}                             & 16 & \emph{(reserved C)}             & 8  & \multicolumn{7}{c}{69 payload} & 77 \\
\bottomrule
\multicolumn{12}{@{}p{\textwidth}@{}}{\vspace{2pt}\footnotesize $^*$Type~5 omits suffix fields to pack two standard callsigns, locator, and SNR into 77 bits ($1+28+28+15+5=77\text{b}$). \quad $^\dagger$16-bit DX hash (Section~\ref{sec:hash}). \quad The ``Total'' column indicates active information bits; formats with Total $< 77$ bits (Types 1, 2, 8, 9, 12) are right-padded with trailing zeros ($77 - \text{Total}$ bits) to complete the 77-bit physical frame boundary prior to CRC-14 computation.}
\end{tabular*}

\vspace{6pt plus 1fill}

% --- Table 4: LQ Varicode Assignments ---
% =============================================================================
% The LQ Digital Mode Family — Reference Implementation (lq_lib)
%
% Author:  Luis Quesada (HB9IPH)
% Web:     https://luisquesada.com
% Portal:  https://lquesada.github.io/lq_lib/
% GitHub:  https://github.com/lquesada/lq_lib
% App:     qFT8 — Portable Amateur Radio for Android (https://qft8.com)
%
% License: MIT License (https://github.com/lquesada/lq_lib/blob/main/LICENSE)
%
% Copyright (c) 2026 Luis Quesada (HB9IPH)
%
% Permission is hereby granted, free of charge, to any person obtaining a copy
% of this software and associated documentation files (the "Software"), to deal
% in the Software without restriction, including without limitation the rights
% to use, copy, modify, merge, publish, distribute, sublicense, and/or sell
% copies of the Software, and to permit persons to whom the Software is
% furnished to do so, subject to the following conditions:
%
% The above copyright notice and this permission notice shall be included in
% all copies or substantial portions of the Software.
%
% THE SOFTWARE IS PROVIDED "AS IS", WITHOUT WARRANTY OF ANY KIND, EXPRESS OR
% IMPLIED, INCLUDING BUT NOT LIMITED TO THE WARRANTIES OF MERCHANTABILITY,
% FITNESS FOR A PARTICULAR PURPOSE AND NONINFRINGEMENT. IN NO EVENT SHALL THE
% AUTHORS OR COPYRIGHT HOLDERS BE LIABLE FOR ANY CLAIM, DAMAGES OR OTHER
% LIABILITY, WHETHER IN AN ACTION OF CONTRACT, TORT OR OTHERWISE, ARISING FROM,
% OUT OF OR IN CONNECTION WITH THE SOFTWARE OR THE USE OR OTHER DEALINGS IN THE
% SOFTWARE.
% =============================================================================
% Auto-generated by tools/generate_varicode.py — do not edit by hand
% Regenerate with: python3 tools/generate_varicode.py --latex -o tables/varicode_table.tex

\caption{LQ Varicode assignments for the 104-character free-text alphabet.}
\label{tab:varicode}
\resizebox{\textwidth}{!}{%
\scriptsize
\setlength{\tabcolsep}{1.2pt}
\begin{tabular}{@{}cl c r | cl c r | cl c r | cl c r@{}}
\toprule
\textbf{Char} & \textbf{Varicode} & \textbf{Len} & \textbf{Freq} & \textbf{Char} & \textbf{Varicode} & \textbf{Len} & \textbf{Freq} & \textbf{Char} & \textbf{Varicode} & \textbf{Len} & \textbf{Freq} & 
\textbf{Char} & \textbf{Varicode} & \textbf{Len} & \textbf{Freq} \\
\midrule
\texttt{{\textvisiblespace}} & \texttt{\szero{}\szero{}1} & 3 & 16.73\% & \texttt{1} & \texttt{\szero{}\szero{}\szero{}\szero{}\szero{}\szero{}1\szero{}11} & 10 & 0.09\% & \texttt{*} & \texttt{\szero{}\szero{}\szero{}\szero{}\szero{}\szero{}\szero{}\szero{}\szero{}\szero{}11\szero{}} & 13 & 0.01\% & \texttt{Å} & \texttt{\szero{}\szero{}\szero{}\szero{}\szero{}\szero{}\szero{}\szero{}\szero{}\szero{}\szero{}\szero{}1\szero{}111} & 17 & $<$0.01\% \\
\texttt{E} & \texttt{\szero{}1\szero{}} & 3 & 10.23\% & \texttt{Q} & \texttt{\szero{}\szero{}\szero{}\szero{}\szero{}\szero{}11\szero{}\szero{}} & 10 & 0.07\% & \texttt{[} & \texttt{\szero{}\szero{}\szero{}\szero{}\szero{}\szero{}\szero{}\szero{}\szero{}\szero{}111\szero{}} & 14 & 0.01\% & \texttt{Æ} & \texttt{\szero{}\szero{}\szero{}\szero{}\szero{}\szero{}\szero{}\szero{}\szero{}\szero{}\szero{}\szero{}11\szero{}\szero{}\szero{}} & 17 & $<$0.01\% \\
\texttt{T} & \texttt{\szero{}11\szero{}} & 4 & 7.44\% & \texttt{!} & \texttt{\szero{}\szero{}\szero{}\szero{}\szero{}\szero{}11\szero{}1} & 10 & 0.07\% & \texttt{]} & \texttt{\szero{}\szero{}\szero{}\szero{}\szero{}\szero{}\szero{}\szero{}\szero{}\szero{}1111} & 14 & 0.01\% & \texttt{Ç} & \texttt{\szero{}\szero{}\szero{}\szero{}\szero{}\szero{}\szero{}\szero{}\szero{}\szero{}\szero{}\szero{}11\szero{}\szero{}1} & 17 & $<$0.01\% \\
\texttt{A} & \texttt{\szero{}111} & 4 & 6.59\% & \texttt{?} & \texttt{\szero{}\szero{}\szero{}\szero{}\szero{}\szero{}111\szero{}} & 10 & 0.07\% & \texttt{$|$} & \texttt{\szero{}\szero{}\szero{}\szero{}\szero{}\szero{}\szero{}\szero{}\szero{}\szero{}\szero{}1\szero{}\szero{}} & 14 & 0.01\% & \texttt{È} & \texttt{\szero{}\szero{}\szero{}\szero{}\szero{}\szero{}\szero{}\szero{}\szero{}\szero{}\szero{}\szero{}11\szero{}1\szero{}} & 17 & $<$0.01\% \\
\texttt{O} & \texttt{1\szero{}\szero{}\szero{}} & 4 & 6.24\% & \texttt{Z} & \texttt{\szero{}\szero{}\szero{}\szero{}\szero{}\szero{}1111} & 10 & 0.06\% & \texttt{/} & \texttt{\szero{}\szero{}\szero{}\szero{}\szero{}\szero{}\szero{}\szero{}\szero{}\szero{}\szero{}1\szero{}1\szero{}} & 15 & $<$0.01\% & \texttt{É} & \texttt{\szero{}\szero{}\szero{}\szero{}\szero{}\szero{}\szero{}\szero{}\szero{}\szero{}\szero{}\szero{}11\szero{}11} & 17 & $<$0.01\% \\
\texttt{N} & \texttt{1\szero{}\szero{}1} & 4 & 5.96\% & \texttt{,} & \texttt{\szero{}\szero{}\szero{}\szero{}\szero{}\szero{}\szero{}1\szero{}\szero{}\szero{}} & 11 & 0.06\% & \texttt{\$} & \texttt{\szero{}\szero{}\szero{}\szero{}\szero{}\szero{}\szero{}\szero{}\szero{}\szero{}\szero{}1\szero{}11\szero{}} & 16 & $<$0.01\% & \texttt{Ê} & \texttt{\szero{}\szero{}\szero{}\szero{}\szero{}\szero{}\szero{}\szero{}\szero{}\szero{}\szero{}\szero{}111\szero{}\szero{}} & 17 & $<$0.01\% \\
\texttt{I} & \texttt{1\szero{}1\szero{}} & 4 & 5.90\% & \texttt{.} & \texttt{\szero{}\szero{}\szero{}\szero{}\szero{}\szero{}\szero{}1\szero{}\szero{}1} & 11 & 0.05\% & \texttt{+} & \texttt{\szero{}\szero{}\szero{}\szero{}\szero{}\szero{}\szero{}\szero{}\szero{}\szero{}\szero{}1\szero{}111} & 16 & $<$0.01\% & \texttt{Ë} & \texttt{\szero{}\szero{}\szero{}\szero{}\szero{}\szero{}\szero{}\szero{}\szero{}\szero{}\szero{}\szero{}111\szero{}1} & 17 & $<$0.01\% \\
\texttt{S} & \texttt{1\szero{}11} & 4 & 5.41\% & \texttt{{\szero}} & \texttt{\szero{}\szero{}\szero{}\szero{}\szero{}\szero{}\szero{}1\szero{}1\szero{}} & 11 & 0.05\% & \texttt{\textbackslash{}n} & \texttt{\szero{}\szero{}\szero{}\szero{}\szero{}\szero{}\szero{}\szero{}\szero{}\szero{}\szero{}11\szero{}\szero{}\szero{}} & 16 & $<$0.01\% & \texttt{Ì} & \texttt{\szero{}\szero{}\szero{}\szero{}\szero{}\szero{}\szero{}\szero{}\szero{}\szero{}\szero{}\szero{}1111\szero{}} & 17 & $<$0.01\% \\
\texttt{R} & \texttt{11\szero{}\szero{}} & 4 & 5.00\% & \texttt{2} & \texttt{\szero{}\szero{}\szero{}\szero{}\szero{}\szero{}\szero{}1\szero{}11} & 11 & 0.05\% & \texttt{\%} & \texttt{\szero{}\szero{}\szero{}\szero{}\szero{}\szero{}\szero{}\szero{}\szero{}\szero{}\szero{}11\szero{}\szero{}1} & 16 & $<$0.01\% & \texttt{Í} & \texttt{\szero{}\szero{}\szero{}\szero{}\szero{}\szero{}\szero{}\szero{}\szero{}\szero{}\szero{}\szero{}11111} & 17 & $<$0.01\% \\
\texttt{H} & \texttt{11\szero{}1} & 4 & 4.76\% & \texttt{8} & \texttt{\szero{}\szero{}\szero{}\szero{}\szero{}\szero{}\szero{}11\szero{}\szero{}} & 11 & 0.04\% & \texttt{\&} & \texttt{\szero{}\szero{}\szero{}\szero{}\szero{}\szero{}\szero{}\szero{}\szero{}\szero{}\szero{}11\szero{}1\szero{}} & 16 & $<$0.01\% & \texttt{Î} & \texttt{\szero{}\szero{}\szero{}\szero{}\szero{}\szero{}\szero{}\szero{}\szero{}\szero{}\szero{}\szero{}\szero{}1\szero{}\szero{}\szero{}} & 17 & $<$0.01\% \\
\texttt{D} & \texttt{111\szero{}\szero{}} & 5 & 3.48\% & \texttt{3} & \texttt{\szero{}\szero{}\szero{}\szero{}\szero{}\szero{}\szero{}11\szero{}1} & 11 & 0.04\% & \texttt{$<$} & \texttt{\szero{}\szero{}\szero{}\szero{}\szero{}\szero{}\szero{}\szero{}\szero{}\szero{}\szero{}11\szero{}11} & 16 & $<$0.01\% & \texttt{Ï} & \texttt{\szero{}\szero{}\szero{}\szero{}\szero{}\szero{}\szero{}\szero{}\szero{}\szero{}\szero{}\szero{}\szero{}1\szero{}\szero{}1} & 17 & $<$0.01\% \\
\texttt{L} & \texttt{111\szero{}1} & 5 & 3.21\% & \texttt{4} & \texttt{\szero{}\szero{}\szero{}\szero{}\szero{}\szero{}\szero{}111\szero{}} & 11 & 0.04\% & \texttt{$>$} & \texttt{\szero{}\szero{}\szero{}\szero{}\szero{}\szero{}\szero{}\szero{}\szero{}\szero{}\szero{}111\szero{}\szero{}\szero{}} & 17 & $<$0.01\% & \texttt{Ð} & \texttt{\szero{}\szero{}\szero{}\szero{}\szero{}\szero{}\szero{}\szero{}\szero{}\szero{}\szero{}\szero{}\szero{}1\szero{}1\szero{}} & 17 & $<$0.01\% \\
\texttt{C} & \texttt{1111\szero{}} & 5 & 2.34\% & \texttt{5} & \texttt{\szero{}\szero{}\szero{}\szero{}\szero{}\szero{}\szero{}1111} & 11 & 0.04\% & \texttt{@} & \texttt{\szero{}\szero{}\szero{}\szero{}\szero{}\szero{}\szero{}\szero{}\szero{}\szero{}\szero{}111\szero{}\szero{}1} & 17 & $<$0.01\% & \texttt{Ñ} & \texttt{\szero{}\szero{}\szero{}\szero{}\szero{}\szero{}\szero{}\szero{}\szero{}\szero{}\szero{}\szero{}\szero{}1\szero{}11} & 17 & $<$0.01\% \\
\texttt{U} & \texttt{11111\szero{}} & 6 & 2.24\% & \texttt{9} & \texttt{\szero{}\szero{}\szero{}\szero{}\szero{}\szero{}\szero{}\szero{}1\szero{}\szero{}} & 11 & 0.03\% & \texttt{\textbackslash{}} & \texttt{\szero{}\szero{}\szero{}\szero{}\szero{}\szero{}\szero{}\szero{}\szero{}\szero{}\szero{}111\szero{}1\szero{}} & 17 & $<$0.01\% & \texttt{Ò} & \texttt{\szero{}\szero{}\szero{}\szero{}\szero{}\szero{}\szero{}\szero{}\szero{}\szero{}\szero{}\szero{}\szero{}11\szero{}\szero{}} & 17 & $<$0.01\% \\
\texttt{M} & \texttt{111111} & 6 & 2.05\% & \texttt{6} & \texttt{\szero{}\szero{}\szero{}\szero{}\szero{}\szero{}\szero{}\szero{}1\szero{}1} & 11 & 0.03\% & \texttt{\^{}} & \texttt{\szero{}\szero{}\szero{}\szero{}\szero{}\szero{}\szero{}\szero{}\szero{}\szero{}\szero{}111\szero{}11} & 17 & $<$0.01\% & \texttt{Ó} & \texttt{\szero{}\szero{}\szero{}\szero{}\szero{}\szero{}\szero{}\szero{}\szero{}\szero{}\szero{}\szero{}\szero{}11\szero{}1} & 17 & $<$0.01\% \\
\texttt{F} & \texttt{\szero{}\szero{}\szero{}1\szero{}\szero{}} & 6 & 1.95\% & \texttt{7} & \texttt{\szero{}\szero{}\szero{}\szero{}\szero{}\szero{}\szero{}\szero{}11\szero{}\szero{}} & 12 & 0.03\% & \texttt{`} & \texttt{\szero{}\szero{}\szero{}\szero{}\szero{}\szero{}\szero{}\szero{}\szero{}\szero{}\szero{}1111\szero{}\szero{}} & 17 & $<$0.01\% & \texttt{Ô} & \texttt{\szero{}\szero{}\szero{}\szero{}\szero{}\szero{}\szero{}\szero{}\szero{}\szero{}\szero{}\szero{}\szero{}111\szero{}} & 17 & $<$0.01\% \\
\texttt{W} & \texttt{\szero{}\szero{}\szero{}1\szero{}1} & 6 & 1.63\% & \texttt{\textquotedbl{}} & \texttt{\szero{}\szero{}\szero{}\szero{}\szero{}\szero{}\szero{}\szero{}11\szero{}1} & 12 & 0.03\% & \texttt{\{} & \texttt{\szero{}\szero{}\szero{}\szero{}\szero{}\szero{}\szero{}\szero{}\szero{}\szero{}\szero{}1111\szero{}1} & 17 & $<$0.01\% & \texttt{Õ} & \texttt{\szero{}\szero{}\szero{}\szero{}\szero{}\szero{}\szero{}\szero{}\szero{}\szero{}\szero{}\szero{}\szero{}1111} & 17 & $<$0.01\% \\
\texttt{P} & \texttt{\szero{}\szero{}\szero{}11\szero{}} & 6 & 1.60\% & \texttt{-} & \texttt{\szero{}\szero{}\szero{}\szero{}\szero{}\szero{}\szero{}\szero{}111\szero{}} & 12 & 0.03\% & \texttt{\}} & \texttt{\szero{}\szero{}\szero{}\szero{}\szero{}\szero{}\szero{}\szero{}\szero{}\szero{}\szero{}11111\szero{}} & 17 & $<$0.01\% & \texttt{Ö} & \texttt{\szero{}\szero{}\szero{}\szero{}\szero{}\szero{}\szero{}\szero{}\szero{}\szero{}\szero{}\szero{}\szero{}\szero{}1\szero{}\szero{}} & 17 & $<$0.01\% \\
\texttt{G} & \texttt{\szero{}\szero{}\szero{}111} & 6 & 1.56\% & \texttt{\textquotesingle{}} & \texttt{\szero{}\szero{}\szero{}\szero{}\szero{}\szero{}\szero{}\szero{}1111} & 12 & 0.02\% & \texttt{\~{}} & \texttt{\szero{}\szero{}\szero{}\szero{}\szero{}\szero{}\szero{}\szero{}\szero{}\szero{}\szero{}111111} & 17 & $<$0.01\% & \texttt{Ø} & \texttt{\szero{}\szero{}\szero{}\szero{}\szero{}\szero{}\szero{}\szero{}\szero{}\szero{}\szero{}\szero{}\szero{}\szero{}1\szero{}1} & 17 & $<$0.01\% \\
\texttt{Y} & \texttt{\szero{}\szero{}\szero{}\szero{}1\szero{}} & 6 & 1.46\% & \texttt{\_} & \texttt{\szero{}\szero{}\szero{}\szero{}\szero{}\szero{}\szero{}\szero{}\szero{}1\szero{}\szero{}} & 12 & 0.02\% & \texttt{¡} & \texttt{\szero{}\szero{}\szero{}\szero{}\szero{}\szero{}\szero{}\szero{}\szero{}\szero{}\szero{}\szero{}1\szero{}\szero{}\szero{}\szero{}} & 17 & $<$0.01\% & \texttt{Ù} & \texttt{\szero{}\szero{}\szero{}\szero{}\szero{}\szero{}\szero{}\szero{}\szero{}\szero{}\szero{}\szero{}\szero{}\szero{}11\szero{}} & 17 & $<$0.01\% \\
\texttt{B} & \texttt{\szero{}\szero{}\szero{}\szero{}11} & 6 & 1.18\% & \texttt{;} & \texttt{\szero{}\szero{}\szero{}\szero{}\szero{}\szero{}\szero{}\szero{}\szero{}1\szero{}1} & 12 & 0.02\% & \texttt{¿} & \texttt{\szero{}\szero{}\szero{}\szero{}\szero{}\szero{}\szero{}\szero{}\szero{}\szero{}\szero{}\szero{}1\szero{}\szero{}\szero{}1} & 17 & $<$0.01\% & \texttt{Ú} & \texttt{\szero{}\szero{}\szero{}\szero{}\szero{}\szero{}\szero{}\szero{}\szero{}\szero{}\szero{}\szero{}\szero{}\szero{}111} & 17 & $<$0.01\% \\
\texttt{V} & \texttt{\szero{}\szero{}\szero{}\szero{}\szero{}1\szero{}} & 7 & 0.85\% & \texttt{:} & \texttt{\szero{}\szero{}\szero{}\szero{}\szero{}\szero{}\szero{}\szero{}\szero{}11\szero{}} & 12 & 0.02\% & \texttt{À} & \texttt{\szero{}\szero{}\szero{}\szero{}\szero{}\szero{}\szero{}\szero{}\szero{}\szero{}\szero{}\szero{}1\szero{}\szero{}1\szero{}} & 17 & $<$0.01\% & \texttt{Û} & \texttt{\szero{}\szero{}\szero{}\szero{}\szero{}\szero{}\szero{}\szero{}\szero{}\szero{}\szero{}\szero{}\szero{}\szero{}\szero{}1\szero{}} & 17 & $<$0.01\% \\
\texttt{K} & \texttt{\szero{}\szero{}\szero{}\szero{}\szero{}11\szero{}} & 8 & 0.53\% & \texttt{=} & \texttt{\szero{}\szero{}\szero{}\szero{}\szero{}\szero{}\szero{}\szero{}\szero{}111\szero{}} & 13 & 0.02\% & \texttt{Á} & \texttt{\szero{}\szero{}\szero{}\szero{}\szero{}\szero{}\szero{}\szero{}\szero{}\szero{}\szero{}\szero{}1\szero{}\szero{}11} & 17 & $<$0.01\% & \texttt{Ü} & \texttt{\szero{}\szero{}\szero{}\szero{}\szero{}\szero{}\szero{}\szero{}\szero{}\szero{}\szero{}\szero{}\szero{}\szero{}\szero{}11} & 17 & $<$0.01\% \\
\texttt{{\ensuremath{\lightning}}} & \texttt{\szero{}\szero{}\szero{}\szero{}\szero{}111} & 8 & 0.34\% & \texttt{(} & \texttt{\szero{}\szero{}\szero{}\szero{}\szero{}\szero{}\szero{}\szero{}\szero{}1111} & 13 & 0.01\% & \texttt{Â} & \texttt{\szero{}\szero{}\szero{}\szero{}\szero{}\szero{}\szero{}\szero{}\szero{}\szero{}\szero{}\szero{}1\szero{}1\szero{}\szero{}} & 17 & $<$0.01\% & \texttt{Ý} & \texttt{\szero{}\szero{}\szero{}\szero{}\szero{}\szero{}\szero{}\szero{}\szero{}\szero{}\szero{}\szero{}\szero{}\szero{}\szero{}\szero{}1} & 17 & $<$0.01\% \\
\texttt{X} & \texttt{\szero{}\szero{}\szero{}\szero{}\szero{}\szero{}1\szero{}\szero{}} & 9 & 0.16\% & \texttt{)} & \texttt{\szero{}\szero{}\szero{}\szero{}\szero{}\szero{}\szero{}\szero{}\szero{}\szero{}1\szero{}\szero{}} & 13 & 0.01\% & \texttt{Ã} & \texttt{\szero{}\szero{}\szero{}\szero{}\szero{}\szero{}\szero{}\szero{}\szero{}\szero{}\szero{}\szero{}1\szero{}1\szero{}1} & 17 & $<$0.01\% & \texttt{Þ} & \texttt{\szero{}\szero{}\szero{}\szero{}\szero{}\szero{}\szero{}\szero{}\szero{}\szero{}\szero{}\szero{}\szero{}\szero{}\szero{}\szero{}\szero{}1} & 18 & $<$0.01\% \\
\texttt{J} & \texttt{\szero{}\szero{}\szero{}\szero{}\szero{}\szero{}1\szero{}1\szero{}} & 10 & 0.10\% & \texttt{\#} & \texttt{\szero{}\szero{}\szero{}\szero{}\szero{}\szero{}\szero{}\szero{}\szero{}\szero{}1\szero{}1} & 13 & 0.01\% & \texttt{Ä} & \texttt{\szero{}\szero{}\szero{}\szero{}\szero{}\szero{}\szero{}\szero{}\szero{}\szero{}\szero{}\szero{}1\szero{}11\szero{}} & 17 & $<$0.01\% & \texttt{[FILL]} & \texttt{\szero{}\szero{}\szero{}\szero{}\szero{}\szero{}\szero{}\szero{}\szero{}\szero{}\szero{}\szero{}\szero{}\szero{}\szero{}\szero{}\szero{}\szero{}} & 18 & $<$0.01\% \\
\bottomrule
\end{tabular}%
}

\end{table*}

\subsection{Protocol Boundary Conditions \& Canonical Precedence}
\label{sec:msg:limitations}

To eliminate ambiguity across station profiles, \LQ{} formalizes eight canonical precedence and boundary rules:
\begin{enumerate}[leftmargin=1.4em, itemsep=1.5pt plus 1pt minus 0.5pt, topsep=2pt, parsep=0pt]
  \item \textbf{Strict Canonical Precedence:} Transmitters \textbf{must} use the most compact format fitting caller and target: for CALL, Type~5 requires both stations standard without suffix; Type~6 applies if either has \texttt{/P} or target is non-standard; Type~7 if caller is non-standard. For replies, Type~8 precedes Type~11 and Type~9 precedes Types~10/12 if both are standard.
  \item \textbf{Deterministic Hash Disambiguation:} In Type~6, target hash matches against unsuffixed standard stations are rejected only if caller also lacks a suffix (Type~5 was mandatory); suffixed callers and Type~7 validly target standard stations. In 1-on-1 contacts, standard stations ignore hashed completions (Types~10, 11) from standard peers.
  \item \textbf{Universal Locator and SNR in Type~6:} Type~6 allocates 4 prefix bits and 24 target hash bits, leaving 49 bits for caller callsign, suffix, 15-bit locator, and 5-bit SNR ($4+24+28+1+15+5=77\text{b}$). Suffixed callers deliver locator and SNR on their initial call without extra steps.
  \item \textbf{Clear-Text Identity in Type~7:} Type~7 encodes the target via 20-bit hash ($H_{20} = H_{24} \gg 4$), the non-standard caller via clear-text 48-bit Base-38 representation ($\le 9$ characters), 1-bit suffix (\texttt{/P}), and 5-bit SNR ($3+20+48+1+5 = 77\text{b}$), delivering unhashed identity directly in Step~2.
  \item \textbf{Single-Target Multi-Station Framing:} When confirming one station in Fox mode (Types~11/12), Target~2 duplicates Target~1 (identical hash, and SNR in Type~11), decoded as a single station. In 1-on-1 contacts, Type~8 strictly supersedes Type~11; single-target Type~11 is reserved for non-standard callers or split F/H schedules.
  \item \textbf{Collision Resistance:} Uniform 24-bit callsign hashes ($2^{24} \approx 16.78\times 10^6$ bins) yield $P_{\text{collision}} \le 0.75\%$ across 500 stations, disambiguated via local callsign caches and canonical precedence.
  \item \textbf{Compound Callsigns:} Stations with callsigns exceeding 9 characters (up to 13 characters, e.g., \callsign{3B9/HB9IPH/P}, Category~D) utilize the 69-bit Base-38 field in CQ frames (Type~4); directed replies use CQ or free text (Type~13).
  \item \textbf{Verification \& Compatibility:} Bit-exact codec traces and golden vectors across all 13 message formats are verified in the reference suite~\cite{lqlib} across AWGN, Rayleigh fading, and clock-drift conditions. Compliant decoders treat unassigned opcodes (Types~14--16) as valid framing boundaries, safely ignoring them.
\end{enumerate}
\vspace*{-4.0pt}

\pagebreak
